\scanpage{33}
\mspage{71}

excellence !). Aussi est-il rare que la raison n'en soit elle-même profondément pervertie, et elle l'est en tous cas toujours peu ou prou.

Il n'y a pas à s'en étonner, quand on se rappelle l'omniprésence des pulsions de désirs dans la psyché, et notamment dans toutes les activités créatrices, tant chez soi-même que chez autrui. Aussi, quand la relation au sexe est faussée, toute la vision de la réalité, et avec elle toutes les relations humaines, sont faussées. Ainsi la raison se trouve sommée, soit de démissionner purement et simplement (c'est là souvent l'attitude du mystique), soit de ``rationaliser'' (!) une vision et des relations aberrantes, par des arguments et des considérations tout aussi aberrants. Même chez des êtres d'une intelligence exquise, quand telle pieuse fiction est en jeu, soudain on se trouve confronté à une débilité telle qu'on croit rêver...

35. La Fuite est le remède - ou la Maison des fous

La double-pensée est un exemple typique et extrême de ``\textit{la Fuite}'' dans la psyché. Nous entendons par là un ensemble cohérent de mécanismes, d'attitudes et de comportements de la psyché, qui inexorablement et de façon permanente nous poussent à nous affirmer (avec une force véhémente dans les couches inconscientes, et comme chose allant de soi dans le conscient) des choses qui sont soit manifestement absurdes ou dénuées en elles-mêmes de tout sens (mais dont le \textit{vrai} sens réside dans l'intention inconsciente de brouiller, de noyer un poisson, ou de faire taire une inquiétude...), soit encore à l'exact opposé de \unsure{.|:}ce qu'au fond nous savons pertinemment, ou qui en représentent une déformation flagrante. (Et de plus dans un aspect si essentiel, le plus souvent, que cette version de la réalité que nous nous présentons à nous-même n'en représente plus qu'une carricature plus ou moins grotesque). \textit{Cet étrange syndrome de la Fuite domaine silencieusement et inexorablement la vie de quasiment tous les hommes, de façon plus ou moins tyrannique et totale suivant le degré de maturité et l'état de fidélité ou d'infidélité de l'âme} - et pourtant, d'un commun et universel accord, il est ignoré par tous ! C'est dire que nous fonctionnons et évoluons, avec les airs désinvoltes de gens très distingué et qui savent de quoi ils parlent, dans une gigantesque \textit{Maison de Fous} à l'échelle de la planète entière, sans autres gardes-chiourme que la commune Folie de tous, et où tous les humains seraient les aliénés et les fous-à-lier et internés à vie.

Quand l'âme pour la première fois s'éveille à la réalité de la Maison de Fous, de cette \textit{Fuite} universelle, obstinée, de tous les instants devant tout ce qui effraie la psyché - elle en reste saisie, comme pétrifiée. Il faudra un bon moment avant qu'elle ait fini de se tater et qu'elle en croie ses yeux et ses saines facultés ! Car elle est seule à voir ce qu'elle voit, et elle ne peut s'empêcher de se dire que c'est elle sûrement qui doit être toquée et que tous ces gens si posés et si sérieux et tous d'accord, au fond, sont sains d'esprit.

Dans l'histoire des hommes sur Terre, le premier à oser en croire le témoignage de ses yeux, en voyant la Maison de Fous, a été \textit{Freud}, au moment où il a découvert l'inconscient. Cela a été (me confirme Mère) la plus grande découverte de l'esprit humain sur terre. Cela a été aussi un des actes essentiels pour marquer le commencement du Temps \deleted{de} \unsure{.}la Fin, et pour préparer le Nouvel Age - l'âge de la Libération. C'est cet \unsure{.}acte, et l'oeuvre de fidélité qui en est née, qui ont joué le rôle le plus essentiel dans l'Oeuvre de Salut et de Libération préparée par la mission du Bouddha et
% archive sheet number: 32
